% #########################################################################
% SCHREIBANLEITUNG KAPITEL 5 -- UND DIE ABGRENZUNG DER KAPITEL 3 BIS 8
% Stand 24.08.2026. Ersetzt die Fassung vom 22.08.2026 vollstaendig.
% Vor dem Schreiben lesen.
% #########################################################################
%
% WAS SICH GEGENUEBER DER VORFASSUNG GEAENDERT HAT
%   1. Kapitel 3 kommt vor. Die Vorfassung begann bei Kapitel 4 und wusste
%      nicht, dass 3.1 die Entscheidungseinheit und 3.2 die Zielgroessen und
%      Referenzstufen bereits festlegt. Damit war die Abgrenzung an ihrem
%      wichtigsten Punkt unvollstaendig.
%   2. Neue Regel fuer die Merkmale, Abschnitt 4 dieses Blocks. Sie loest
%      das Unbehagen auf, die Merkmale erst in Kapitel 4 zu beschreiben und
%      danach in Kapitel 5 zu berechnen.
%   3. Drei falsche Angaben der Vorfassung sind hier nicht mehr enthalten;
%      sie stehen zur Sicherheit unter Abschnitt 10 als "nicht wieder
%      hereinschreiben".
%   4. Verweise laufen ueber ABSCHNITTSNAMEN, nicht ueber Nummern. Solange
%      die Entscheidung aus Abschnitt 6 offen ist, waeren Nummern zur
%      Haelfte falsch.
%
% -------------------------------------------------------------------------
% 1. DER ROTE FADEN IN DREI SAETZEN
% -------------------------------------------------------------------------
% Kapitel 3  Was soll entschieden werden? Anwendungsfall, Entscheidungs-
%            einheit, Zielgroessen, Erfolgsmassstab. Fachlich, ohne eine
%            einzige Zahl aus den Daten.
% Kapitel 4  Was liefern die Quellen, und was faellt daran auf?
%            BEFUND, ohne ein einziges "deshalb".
% Kapitel 5  Was wurde damit gemacht, warum, und woran sieht man, dass es
%            stimmt? EINGRIFF und NACHWEIS, ohne eine einzige
%            Verteilungskennzahl neu herzuleiten.
%
% Kapitel 5 ist das einzige Kapitel, das den Nachweis mitliefert: die 19
% automatisierten Pruefungen an den erzeugten Dateien. Genau darin liegt
% seine Rolle. Der Merksatz aus dem Gutachten -- Probleme zu ERKENNEN
% reicht nicht, bewertet wird, ob sie methodisch GELOEST wurden -- ist die
% Naht zwischen Kapitel 4 und Kapitel 5.
%
% MERKSATZ: Ein Befund gehoert einmal genannt und danach nur noch zitiert.
% Wer ihn beim zweiten Mal erneut herleitet, schreibt das Kapitel doppelt.
%
% -------------------------------------------------------------------------
% 2. DIE TRENNUNG IN EINEM SATZ JE KAPITEL
% -------------------------------------------------------------------------
%   Kap. 3  Business Understanding  AUFTRAG   Was ist die Frage, woran misst
%                                             sie sich?  -- FERTIG
%   Kap. 4  Data Understanding      BEFUND    Was ist da, was faellt auf?
%                                             -- FERTIG, nicht mehr anfassen
%   Kap. 5  Data Preparation        EINGRIFF  Was wurde damit gemacht, warum,
%                                             und was beweist es?
%   Kap. 6  Modelling               VERFAHREN Was wurde angewandt, warum das?
%   Kap. 7  Evaluation              ERGEBNIS  Wie gut war es?
%   Kap. 8  Diskussion              DEUTUNG   Was heisst das, was gilt nicht?
%
% VIER TESTFRAGEN, wenn unklar ist, wohin ein Absatz gehoert:
%   0. Gaelte der Satz auch ohne diese Daten, allein aus dem Anwendungsfall?
%                                                      -> ja: Kapitel 3
%   1. Gaelte der Satz auch, wenn die Daten gar nicht verarbeitet wuerden?
%                                                      -> ja: Kapitel 4
%   2. Veraendert dieser Schritt die Datei auf der Platte?
%                                                      -> ja: Kapitel 5
%   3. Passiert das je Fold, innerhalb der sklearn-Pipeline?
%                                                      -> ja: Kapitel 6
%
% -------------------------------------------------------------------------
% 3. WAS KAPITEL 3 UND 4 BEREITS VERBRAUCHT HABEN
% -------------------------------------------------------------------------
% Links steht, was schon dasteht. Rechts, was in Kapitel 5 dazugehoert.
% Links darf NUR NOCH ZITIERT werden -- mit Verweis, ohne Herleitung, ohne
% die Zahl ein zweites Mal zu erklaeren.
%
% DIESE TABELLE IST DIE DRAMATURGIE DES KAPITELS, nicht nur eine Notiz.
% Jeder Abschnitt von Kapitel 5 beginnt mit dem Befund, den er beantwortet
% -- ein Halbsatz, kein Absatz. Beispiel:
%   "Von den 41 Analysis Neighborhoods aus Abschnitt 4.1 gehen 35 in die
%    Analyse ein."
% Dann kommt der Eingriff. Damit hat der Leser durchgehend das Gefuehl,
% dass hier erledigt wird, was vorher offen blieb.
%
%   Steht in Kapitel 3 oder 4              Eingriff in Kapitel 5
%   ------------------------------------   -------------------------------
%   3.1 Entscheidungseinheit               Uebernahme als Analyseeinheit,
%       Stadtteil x Monat, fachlich         Ausschluss der beiden
%       begruendet                          Alternativen (Einzeleinsatz,
%                                           Tract x Jahr). NICHT als neue
%                                           Wahl darstellen -- gewaehlt
%                                           wurde in 3.1
%   3.2 drei Zielgroessen, zwei des        Konstruktion der Zielgroessen
%       Mengen-, eine des Struktur-         (Abschnitt Merkmale)
%       strangs
%   3.2 die Referenzstufen als             Stufe 1 und Stufe 2 samt Werten
%       Erfolgsmassstab                     (Abschnitt Analyserahmen)
%   4.1 sieben Quellen, drei Raumbezuege   die drei Joins samt Trefferquoten
%       und zwei Zeitraster fallen          (Abschnitt Zusammenfuehrung)
%       auseinander
%   4.1 41 Analysis Neighborhoods          Vollzug des Ausschlusses auf 35.
%   4.2 nennt alle SECHS ausgeschlossenen   Die sechs Namen stehen schon in
%       Stadtteile namentlich               4.2 -- hier nur der Vollzug und
%                                           der Unterschied zwischen
%                                           zeilenweisem dropna und dem
%                                           Ausschluss ganzer Einheiten
%   4.2 269 doppelte Einsatznummern         Entfernung, Zahl nicht neu
%       (0,04 %)                            herleiten
%   4.2 Antwortzeit durchgehend             Fenster 0-60 min als
%       plausibel                           Plausibilitaetspruefung, OHNE
%                                           Entfallquote. Es entfaellt
%                                           keine einzige Zeile. Dazu der
%                                           Satz, dass die Grenze GESETZT
%                                           ist und nicht aus den Daten folgt
%   4.2 10,9 % der Parzellen ohne Baujahr   Nenner yrbuilt_count statt
%                                           parcel_count
%   4.2 ACS ist Periodenschaetzung mit      Regel acs_jahr kleinergleich
%       Publikationsversatz                 Einsatzjahr minus 1
%   4.2 ACS 2009 ohne Bildungsangaben       Analysebeginn 2015-01 als
%                                           KONSEQUENZ, nicht als Setzung
%   4.2 Tract-Grenzen aendern sich,         die drei Trefferquoten
%       Crosswalk beruht auf Zensus 2020    63,1 / 79,7 / 99,2 %
%   4.2 Klassenverteilung                   argmax-Bildung der vier Klassen
%       79,0 / 16,3 / 3,1 / 1,5 %
%   4.2 Dispersionsindex 62,8,              Poisson-GLM als Stufe-2-Baseline
%       rechtsschiefe Zielgroesse
%   4.2 92,5 % Zwischenvarianz              Stadtteil-Split -- steht als
%                                           Spalte IN der Datei
%   4.2 zeitliche Aufloesung 128,3 / 4 / 1  NICHTS in Kapitel 5. Der Befund
%                                           wird in 7.4 und 8.3 zitiert
%
% -------------------------------------------------------------------------
% 4. DIE MERKMALSREGEL -- die wichtigste Neuerung dieses Blocks
% -------------------------------------------------------------------------
% Kapitel 4 beschreibt die zwoelf Modellmerkmale, Kapitel 5 berechnet sie.
% Das ist die richtige Reihenfolge -- man kann nichts explorieren, was man
% nicht konstruiert hat, und keine Konstruktion begruenden, die man nicht
% exploriert hat. CRISP-DM zeichnet zwischen Data Understanding und Data
% Preparation deshalb Pfeile in beide Richtungen.
%
% Falsch wird es erst, wenn beide Kapitel auf DERSELBEN FLUGHOEHE
% beschreiben. Dann liest sich das zweite wie eine Wiederholung. Die Regel:
%
%   KAPITEL 4 NENNT KEINE RECHENVORSCHRIFT.
%   KAPITEL 5 NENNT KEINE VERTEILUNGSKENNZAHL.
%
%   Merkmal                       Kapitel 4          Kapitel 5
%   ---------------------------   ----------------   --------------------
%   armutsquote_pct               Korrelation        Zaehler und Nenner,
%                                 0,465, vier Werte  Nenner kleinergleich 0
%                                 je Stadtteil       ergibt NaN
%   log_kriminalitaetsindex       128,3 Werte je     Location Quotient,
%                                 Stadtteil, das     Fenster endend im
%                                 einzige monatlich  Vormonat, logarithmiert
%                                 variierende
%                                 Merkmal, rho 0,589
%   anteil_altbau_vor_1940_pct    ein Wert je        Nenner ist
%                                 Stadtteil,         yrbuilt_count, nicht
%                                 Zwischenvarianz    parcel_count
%                                 1,000
%
% Rechts steht nie eine Zahl ueber die Verteilung, links nie eine Formel.
% Kein Satz erscheint zweimal, und der Leser denkt beim zweiten Durchgang
% nicht "das hatten wir schon", sondern "daher kommt das also".
%
% FOLGE FUER DIE MERKMALSTABELLE IN 4.1 (offener Punkt P1): Sie darf
% Skalenniveau, Einheit und Wertebereich fuehren -- aber KEINE Spalte
% "Berechnung".
%
% -------------------------------------------------------------------------
% 5. WAS NICHT IN KAPITEL 5 GEHOERT
% -------------------------------------------------------------------------
%   Linearitaetspruefung, Residuen, RESET-Test   -> 6.2, prueft ein
%                                                   VERFAHREN, nicht die Daten
%   VIF / Multikollinearitaet                    -> 6.2, betrifft Ridge
%   StandardScaler, log(1+y), Label-Encoder      -> 6, laeuft je Fold
%   class_weight "balanced"                      -> 6, Modellparameter
%   Hyperparameter-Suchraeume, Tuning-Budget     -> 6
%   Fold-Ergebnisse der Verfahren                -> 7
%   Trainings- und Inferenzzeiten                -> 7
%   Vergleich der Verfahren, Rangfolge oder      -> 8
%     deren Verweigerung bei ueberlappender Streuung
%   n_eff rund 38 (ICC 0,926), Extrapolation     -> 8, Limitation
%
% WAS SEHR WOHL HIERHER GEHOERT: die BASELINES samt ihren Werten
% (R2 0,542 / RMSE 33,98 usw.). Auflage Schroeter vom 27.07.2026 -- sie
% gehoeren in die Data Preparation, nicht nach 6 oder 7. CRISP-DM stuetzt
% das: unter "Select data" steht als Aktivitaet ausdruecklich "Consider use
% of sampling techniques (e.g. ... splitting test and training datasets)".
% Der Stadtteil-Split steht also nicht nur wegen der Auflage hier, sondern
% weil das Vorgehensmodell ihn selbst in dieser Phase fuehrt. Ein Satz dazu
% im Abschnitt Analyserahmen nimmt der Frage "warum steht Ihre Validierung
% nicht in Kapitel 6?" die Spitze.
%
% NOTIZ ZUR EXPOSE-ANKUENDIGUNG:
% Das Expose kuendigte an, "bei ausgepraegten Ungleichgewichten werden
% seltene Kategorien zusammengefasst oder die Klassifikation auf zwei
% Klassen vereinfacht, etwa Brand vs. Nicht-Brand". Das ist NICHT
% eingetreten: Die Zielgroesse hat vier Klassen, die Schieflage wird ueber
% class_weight "balanced" und Macro-F1 behandelt, nicht ueber
% Zusammenlegung. In Kapitel 5 steht davon nur die argmax-Bildung; die
% Behandlung gehoert nach 6, die Abweichung vom Expose nach 8.
%
% -------------------------------------------------------------------------
% 6. GLIEDERUNG VON KAPITEL 5 -- eine Entscheidung ist offen
% -------------------------------------------------------------------------
% CRISP-DM fuehrt fuer die Phase Data Preparation fuenf Aufgaben
% (Chapman et al. 2000, S. 23-26, im Referenzmodell nachgeschlagen):
% 3.1 Select data, 3.2 Clean data, 3.3 Construct data, 3.4 Integrate data,
% 3.5 Format data. Die Ueberschriften treffen sie fast eins zu eins:
%
%   Datenauswahl und Bereinigung        Select data + Clean data
%   Zusammenfuehrung der Datenquellen   Integrate data
%   Konstruktion der Merkmale und
%     Zielgroessen                      Construct data
%   Analyserahmen und Baselines         Select data (Splitting) + Format data
%
% Damit ist die Kolloquiumsfrage "warum ist Ihr Kapitel 5 so geschnitten?"
% genauso beantwortet wie bei Kapitel 4: die Gliederung folgt dem
% Vorgehensmodell. Eine Fussnote, \footcite[23-26]{Chapman2000}.
% NEBENBEI ZU BENENNEN: Die Reihenfolge weicht ab -- integriert wird VOR
% dem Konstruieren, weil die Pipeline es so tut (s1 joint auf Einsatzebene,
% s2 aggregiert und konstruiert danach). Ein Halbsatz genuegt; ungesagt
% liest es sich als Schludrigkeit statt als bewusste Ordnung.
%
% OFFEN: WO STEHT DIE ANALYSEEINHEIT?
% Die 13 Zeilen Fliesstext, die heute zwischen \section{Data Preparation}
% und der ersten \subsection stehen, verstossen gegen die Blindabsatz-Regel
% (Abschnitt 8) und muessen ohnehin weg. Drei Varianten:
%
%   A  eigener Abschnitt "Analyseeinheit" davor, alle uebrigen um eins
%      verschoben. Die Entscheidung steht im Inhaltsverzeichnis -- das
%      zaehlt fuer das Bewertungskriterium "kritische Reflexion im Verlauf"
%      (Gutachten 2,7). Preis: fuenf Unterabschnitte, und alle
%      Kapitelverweise im Code verschieben sich um +1.
%   B  Abschnitt "Analyseeinheit und Datenauswahl", Bereinigung faellt aus
%      dem Titel und bleibt im Inhalt. Vier Unterabschnitte, Entscheidung
%      trotzdem im Verzeichnis sichtbar.
%   C  Alles unter "Datenauswahl und Bereinigung", Analyseeinheit als
%      erster Absatz. Nummerierung bleibt 5.1 bis 5.4, keine Codeverweise
%      zu aendern. Preis: die Entscheidung ist im Verzeichnis unsichtbar.
%
% Seit klar ist, dass 3.1 die Entscheidungseinheit BEREITS begruendet,
% spricht mehr fuer C als vorher: Kapitel 5 trifft hier keine Wahl, es
% setzt eine getroffene Wahl um. Bei C beginnt der Abschnitt mit der
% Analyseeinheit -- der erste Satz des heutigen Fliesstextes
% ("Die Data-Preparation-Phase ueberfuehrt ...") faellt ersatzlos weg, er
% ist eine Kapitelankuendigung und damit ein Blindabsatz in Satzform.
%
% NOCH NICHT ENTSCHIEDEN. Bis dahin: Abschnitte im Text ueber \ref
% ansprechen, nicht ueber ausgeschriebene Nummern.
%
% EMPFEHLUNG ZU EINER UEBERSCHRIFT: "Zusammenfuehrung der VIER
% Datenquellen" -> die Zahl streichen. Kapitel 4 fuehrt in Tabelle 4.1
% sieben Quellen; die Vier meint ACS, Kriminalitaet, Land Use und
% Geometrie, also ohne die Einsatzdaten und ohne den Crosswalk. Beide
% Zaehlungen sind fuer sich richtig, nebeneinander sehen sie nach einem
% Fehler aus -- und die Ueberschrift ist die exponierteste Stelle dafuer.
%
% ENTSCHIEDEN AM 24.08.2026: "Baselines" bleibt stehen. Der Begriff ist in
% Kapitel 7 und im Code ("Stufe-2-Baseline") durchgezogen; ein Wechsel auf
% "Referenzmodelle" nur hier waere schlechter als gar keiner. Falls
% Anglizismen abgebaut werden, dann gesammelt im Formalia-Durchgang M6.
%
% -------------------------------------------------------------------------
% 7. ZWEI SONDERFAELLE
% -------------------------------------------------------------------------
%  (a) 4.2 EDA gegen 6.2 Eignungspruefung. Beide schauen Verteilungen an.
%      Die Regel: In 4.2 steht, WIE DIE DATEN BESCHAFFEN SIND ("die
%      Zielgroesse ist rechtsschief, Dispersionsindex 62,8"). In 6.2 steht,
%      OB EIN VERFAHREN DAZU PASST ("weil sie rechtsschief ist, wird Ridge
%      auf log(1+y) geschaetzt; der RESET-Test verwirft die lineare
%      Spezifikation"). Jede Abbildung erscheint GENAU EINMAL.
%      ACHTUNG: A16 und A17 sind in 4.2 verbraucht -- 6.2 braucht eigene
%      Abbildungen oder verweist auf Kapitel 4.
%
%  (b) 2.3 Grundlagen der Verfahren gegen 6. Kapitel 2 erklaert, WIE ein
%      Verfahren funktioniert -- allgemein, lehrbuchhaft, ohne diesen
%      Datensatz. Kapitel 6 sagt, WARUM es hier eingesetzt wird und MIT
%      WELCHEN Einstellungen. Wer in 6 erklaeren muss, wie Boosting
%      arbeitet, hat es in 2.3 versaeumt.
%      Dasselbe gilt fuer den Verfuegbarkeitsgrundsatz: Dass ein Praediktor
%      zum Prognosezeitpunkt verfuegbar gewesen sein muss, ist Lehrbuch-
%      wissen und gehoert nach Kapitel 2. In Kapitel 5 steht die ANWENDUNG
%      -- welche Spalten, welcher Versatz, welcher Split -- mit einem
%      Verweis, ohne den Grundsatz erneut herzuleiten.
%
% -------------------------------------------------------------------------
% 8. BLINDABSAETZE -- GILT FUER JEDE GLIEDERUNGSEBENE DER ARBEIT
% -------------------------------------------------------------------------
% Woertlich aus der Mitschrift vom 27.07.2026:
%   "Blindabsaetze: Zwischen 1 und 1.1 steht kein Inhalt -> redundanz
%    vermeiden"
% und aus der Mitschrift vom 10.08.2026:
%   "Keine zwischenfazits (keine reudundanz)! Blindabsaetze nutzen!"
%
% Ein BLINDABSATZ ist die BEWUSST LEERE Stelle zwischen einer Ueberschrift
% und ihrer ersten Unterueberschrift. "Blindabsaetze nutzen" heisst: diese
% Stelle LEER LASSEN.
%   Zwischen \section und dem ersten \subsection steht KEIN Text.
%   Zwischen \subsection und dem ersten \subsubsection steht KEIN Text.
% BEGRUENDUNG (Schroeter): Ein Absatz an dieser Stelle kuendigt an, was
% gleich ohnehin kommt -- ein vorweggenommenes Zwischenfazit. Der Inhalt
% gehoert in den ersten Unterabschnitt, wo er ohne Doppelung steht.
%
% Die Regel gilt auch fuer EINZELNE SAETZE. Ein Satz, der beschreibt, was
% das Kapitel gleich tun wird, ist ein Blindabsatz in kurz.
%
% -------------------------------------------------------------------------
% 9. CODEAUSSCHNITTE
% -------------------------------------------------------------------------
% Sieben Ausschnitte, fertig als lstlisting in docs/kapitel5_listings.tex,
% Auswahl und Overleaf-Einbindung in docs/Kapitel5_Codeausschnitte.md.
% Zwei Regeln daraus:
%   - Jede Zeile steht ZEICHENGLEICH so in der .py-Datei. Gekuerzt wird nur
%     durch Weglassen ganzer Zeilen, nie durch Umformulieren. Soll ein
%     Kommentar anders lauten, wird er ZUERST im Code geaendert.
%   - Der Ausschnitt zeigt eine ENTSCHEIDUNG. Wer ihn mit "das laedt eine
%     Datei" zusammenfassen kann, hat ihn falsch gewaehlt.
% Seitenbudget: Code zaehlt zu den 40-60 Inhaltsseiten, Abbildungen nicht.
% Der Anhang zaehlt nicht mit und soll moeglichst entfallen (22.08.2026) --
% es gibt also keinen Fluchtweg fuer zu lange Ausschnitte.
%
% -------------------------------------------------------------------------
% 10. ARBEITSWEISE -- und was nicht wieder hereinschreiben
% -------------------------------------------------------------------------
% Alle ZAHLEN kommen aus docs/03_STAND.md -- dort stehen sie, und nur dort.
% Nicht aus aelteren Fassungen abschreiben. Die Pipeline wurde am
% 28.07.2026 vollstaendig umgebaut.
% Reihenfolge des Kapitels folgt der Pipeline:
%   prep/s1_daten.py -> prep/s2_datensaetze.py -> vorpruefung/v1_baselines.py
%
% DREI ANGABEN DER VORFASSUNG WAREN FALSCH. Sie stehen teils noch in den
% INHALT-Bloecken unter den einzelnen \subsection und sind dort zu ziehen:
%
%   (a) "Ausrueckzeit auf 0-60 min begrenzt (~1,7 % entfallen)".
%       Es heisst ANTWORTZEIT (arrival_dttm minus alarm_dttm, Spalte
%       response_time_min), und es entfaellt KEINE EINZIGE ZEILE:
%       720.258 minus 269 Dubletten ergibt exakt die 719.989 aus
%       03_STAND.md. Die 1,7 % stammen aus einem Stand vor dem
%       Downloadfilter "arrival_dttm IS NOT NULL".
%
%   (b) "LISTING 1: safe_ratio". Eine Funktion dieses Namens gibt es nicht.
%       Der Schutz gegen Division durch null sitzt in
%       s1_daten.berechne_quoten als n.where(n > 0, np.nan).
%
%   (c) "Reihenfolge folgt der Pipeline: ... -> s3_baselines.py".
%       prep/s3_baselines.py existiert nicht. Die Baselines stehen in
%       vorpruefung/v1_baselines.py -- in einem eigenen Ordner VOR der
%       Modellierung. Genau so ist Schroeters Auflage im Code umgesetzt,
%       und genau so gehoert es in den Text.
%
% NACHZUZIEHEN IM CODE (Kommentare, kein Verhaltensrisiko, kein Neulauf):
%   tests/test_aufbereitung.py:127  "NegBin-Offset" -> "Poisson-Offset"
%   tests/test_aufbereitung.py:5-6  Bezug auf den Anhang streichen
%   prep/config.py:76-77            ergaenzen, dass 0-60 min eine GESETZTE
%                                   Grenze ist und nicht aus den Daten folgt
%   fuenf Kapitelverweise           nur falls Variante A oder B gewaehlt
%     (prep/config.py:114, prep/s1_daten.py:451,
%      prep/s2_datensaetze.py:39 und :110, tests/test_aufbereitung.py:127)
% #########################################################################
